\subsection{豪斯多夫级数中的代入}

由于 K 是含有 Q 的域，豪斯多夫级数可看作系数在 K 中的李形式幂级数。因此，若 g 是满足 \(g = \bigcup_{\alpha > 0} g_{\alpha}\) 的完备豪斯多夫滤李代数，则对于 g 中的 a、b，可以在 H 中以 a、b 代替 U、V（参见第 3 号和 §2 第 5 号的注）。

特别地，令 A 为完备豪斯多夫滤含幺结合代数。对 \(m = \bigcup_{\alpha > 0} A_{\alpha}\)，记 \(m_{\alpha} = A_{\alpha} \cap m\) 和 \(\alpha \in R\)；因此当 \(m_{\alpha} = A_{\alpha}\) 时 \(\alpha > 0\)，而当 \(m_{\alpha} = m\) 时 \(\alpha \leqslant 0\)。赋予 m 括号 \([a, b] = ab - ba\) 后，m 是完备豪斯多夫滤李代数，因而可以应用上述结果。用此记号，得到如下结果，它补充了第 1 号的命题 1。

命题 3。若 \(a \in \mathfrak{m}\)、\(b \in \mathfrak{m}\)，则 \(\exp \mathrm{H}(a, b) = \exp a \cdot \exp b\)。

令 \(a, b\) 属于 \(\mathfrak{m}\)；存在 \(\alpha > 0\)，使得 \(a \in \mathrm{A}_{\alpha}\) 且 \(b \in \mathrm{A}_{\alpha}\)。于是存在连续同态 \(\theta\)，它从 Magnus 代数 \(\hat{\mathbf{A}}(\{\mathbf{U},\mathbf{V}\})\) 到 A，将 U 映到 \(a\)，将 V 映到 \(b\)（ \(\S 5\) ，第 1 号，命题 1）。

将 \(\theta\) 限制在 \(\hat{\mathbf{L}}(\{\mathrm{U},\mathrm{V}\})\) 上，得到从 \(\mathbf{L}(\{\mathrm{U},\mathrm{V}\})\) 的李代数到 \(\mathfrak{m}\) 的连续同态，它将 \(\mathbf{U}\)（分别为 V）映到 \(a\)（分别为 \(b\)）。因此由第 3 号的公式 (6)，\(\theta (\mathrm{H}) = \mathrm{H}(a,b)\)。只需将连续同态 \(\theta\) 施加于关系式两边

\[
\exp \mathrm{H} (\mathrm{U}, \mathrm{V}) = \exp \mathrm{U}. \exp \mathrm{V}
\]

并注意第 1 号的注 3。

注 1。若 \(a\) 与 \(b\) 可交换，则当 \(\mathrm{H}_{r,s}(a,b) = 0\) 时，\(r + s \geqslant 2\)，因为每个次数 \(\geqslant 2\) 的齐次李多项式在 \((a,b)\) 处均为零。于是 \(\mathrm{H}(a,b) = a + b\)，命题 3 恢复了公式

\[
\exp (a + b) = \exp a. \exp b.
\]

命题 4。设 \(\mathfrak{g}\) 是满足 \(\mathfrak{g} = \bigcup_{\alpha >0} \mathfrak{g}_{\alpha}\) 的完备豪斯多夫滤李代数。映射 \((a, b) \mapsto \mathrm{H}(a, b)\) 是 \(\mathfrak{g}\) 上与 \(\mathfrak{g}\) 上的拓扑相容的群律；在该拓扑下，\(0\) 是单位元，且 \(-a\) 是 \(a\) 的逆元，其中对所有 \(a \in \mathfrak{g}\) 均成立。

从 \((a, b) \mapsto \mathrm{H}(a, b)\) 到 \(\mathfrak{g} \times \mathfrak{g}\) 的映射 \(\mathfrak{g}\) 是连续的（第 3 号）；由于映射 \(a \mapsto -a\) 显然连续，只需证明下列关系

\begin{enumerate}
\def\labelenumi{(\arabic{enumi})}
\setcounter{enumi}{17}
\tightlist
\item
\end{enumerate}

\[
\mathrm{H} (\mathrm{H} (a, b), c) = \mathrm{H} (a, \mathrm{H} (b, c))\tag{18}
\]

\[
\mathrm{H} (a, - a) = 0\tag{19}
\]

\[
\mathrm{H} (a, 0) = \mathrm{H} (0, a) = a\tag{20}
\]

对于 \(a, b, c\) 中的 \(\mathfrak{g}\)。由第 3 号公式 (7)，只需在 \(a, b, c\) 为三个不定元且 \(\mathfrak{g} = \hat{\mathbf{L}}\langle \{a, b, c\}\rangle\) 时证明这些公式。现在，指数映射限制在 \(\hat{\mathbf{L}}\langle \{a, b, c\}\rangle\) 上，是到 Magnus 代数 \(\hat{\mathbf{A}}\langle \{a, b, c\}\rangle\) 中的一个嵌入，并且由命题 3：

\[
\exp \mathrm{H} (\mathrm{H} (a, b), c) = \exp \mathrm{H} (a, b). \exp c = \exp a. \exp b. \exp c
\]

\[
\exp \mathrm{H} (a, \mathrm{H} (b, c)) = \exp a. \exp \mathrm{H} (b, c) = \exp a. \exp b. \exp c
\]

\[
\exp \mathrm{H} (a, - a) = \exp a. \exp (- a) = \exp (a - a) = \exp 0
\]

\[
\exp \mathrm{H} (a, 0) = \exp a. \exp 0 = \exp a
\]

\[
\exp \mathrm{H} (0, a) = \exp 0. \exp a = \exp a.
\]

这就证明了关系 (18) 至 (20)。

注。（2）取 g 为李代数 \(\hat{\mathbf{L}}(\mathbf{X})\)。上述命题中引入的群律与第 2 号定义的运算一致。换言之，

\[
a \uplus b = \mathrm{H} (a, b) \quad \text { for   } a, b \text {   in   } \hat {\mathrm{L}} (\mathrm{X});\tag{21}
\]

因此，豪斯多夫群律由豪斯多夫级数给出。

\begin{enumerate}
\def\labelenumi{(\arabic{enumi})}
\setcounter{enumi}{2}
\tightlist
\item
  设 \(\mathfrak{g}\) 是一个带有由下中心列定义的整滤子的李代数 \((\mathcal{C}^m\mathfrak{g})\)。假设存在 \(m \geqslant 1\)，使得 \(\mathcal{C}^m\mathfrak{g} = \{0\}\)。赋予其由滤子 \((\mathcal{C}^m\mathfrak{g})_{n \geqslant 1}\) 导出的拓扑后，\(\mathfrak{g}\) 是豪斯多夫的、完备的，甚至是离散的。于是有 \(\mathrm{P}(a_1, \ldots, a_r) = 0\)，其中 \(a_1, \ldots, a_r\) 属于 \(\mathfrak{g}\)，且这里的 \(\mathrm{P}\) 是次数 \(\geqslant m\) 的任意齐次李多项式；特别地，\(\mathrm{H}_{r,s}(a, b) = 0\) 当 \(r + s \geqslant m\) 时，且级数 \(\mathrm{H}(a, b) = \sum_{r,s} \mathrm{H}_{r,s}(a, b)\) 只有有限个非零项。于是，映射 \((a, b) \mapsto \mathrm{H}(a, b)\) 在 \(\mathfrak{g}\) 上是一个群律，并且是多项式映射（§ 2，第 4 号）。
\end{enumerate}

命题 5。令 \(\mathbf{K}_{r,s}\) 为 \(\mathrm{H}(\mathrm{U} + \mathrm{V}, -\mathrm{U})\) 的多重次数为 \((r,s)\) 的分量。则

\[
\mathrm{K} _ {n, 1} (\mathrm{U}, \mathrm{V}) = \frac {1}{(n + 1) !} (\operatorname{ad} \mathrm{U}) ^ {n} (\mathrm{V}) \quad for n \geqslant 0.
\]

记 \(\mathrm{K}(\mathrm{U},\mathrm{V}) = \mathrm{H}(\mathrm{U} + \mathrm{V}, - \mathrm{U}),\mathrm{K}_1(\mathrm{U},\mathrm{V}) = \sum_{n\geqslant 0}\mathrm{K}_{n,1}(\mathrm{U},\mathrm{V}).\) 记 L（分别为 R）为 \(\hat{\mathbf{A}} (\{\mathbf{U},\mathbf{V}\})\) 上由 U 左乘（分别右乘）的映射

 可写为

\[
\begin{array}{r l} e ^ {\mathrm{U}} \mathrm{V} e ^ {- \mathrm{U}} & = \sum_ {p, q} \frac {\mathrm{U} ^ {p}}{p !} \mathrm{V} \frac {(- \mathrm{U}) ^ {q}}{q !} \\ & = \sum_ {n \geqslant 0} \frac {1}{n !} \left(\sum_ {p + q = n} \frac {n !}{p ! q !} (\mathrm{L} ^ {p} (- \mathrm{R}) ^ {q}). \mathrm{V}\right) \\ & = \sum_ {n \geqslant 0} \frac {1}{n !} (\mathrm{L} - \mathrm{R}) ^ {n}. \mathrm{V} \end{array}
\]

从而

\[
e ^ {\mathrm{U}} \mathrm{V} e ^ {- \mathrm{U}} = \sum_ {n \geqslant 0} \frac {1}{n !} (\mathrm{adU}) ^ {n} \mathrm{V}.\tag{22}
\]

现在模去 \(\sum_{m\geqslant 0}\sum_{n\geqslant 2}\mathrm{A}^{m,n}(\{\mathrm{U},\mathrm{V}\})\) 的理想 \(\mathrm{A}(\{\mathrm{U},\mathrm{V}\})\) 进行计算。对所有 \(n\geqslant 1\)，

\[
(\mathrm{U} + \mathrm{V}) ^ {n} \equiv \mathrm{U} ^ {n} + \sum_ {i = 1} ^ {n - 1} \mathrm{U} ^ {i} \mathrm{V} \mathrm{U} ^ {n - 1 - i}
\]

从而

\[
\begin{array}{r l} (\mathrm{adU}) (\mathrm{U} + \mathrm{V}) ^ {n} & \equiv ((\mathrm{L} - \mathrm{R}) \sum_ {i = 1} ^ {n - 1} \mathrm{L} ^ {i} \mathrm{R} ^ {n - i}). \mathrm{V} \\ & \equiv (\mathrm{L} ^ {n} - \mathrm{R} ^ {n}). \mathrm{V} \\ & \equiv \mathrm{U} ^ {n} \mathrm{V} - \mathrm{VU} ^ {n}. \end{array}
\]

因此

\[
(\mathrm{ad} \mathrm{U}). \mathrm{e} ^ {\mathrm{U} + \mathrm{V}} \equiv e ^ {\mathrm{U}} \mathrm{V} - \mathrm{V} e ^ {\mathrm{U}}\tag{23}
\]

对 n 求和。

另一方面，\(\mathrm{K}_1(\mathrm{U},\mathrm{V})\equiv \mathrm{K}(\mathrm{U},\mathrm{V})\)，且 \(e^{\mathrm{K}_1(\mathrm{U},\mathrm{V})}\equiv 1 + \mathrm{K}_1(\mathrm{U},\mathrm{V})\)，从而

\[
\mathrm{K} _ {1} \equiv e ^ {\mathrm{K}} - 1 \equiv e ^ {\mathrm{U} + \mathrm{V}} e ^ {- \mathrm{U}} - 1
\]

由命题 3 得到

\[
\begin{array}{r l r} (\mathrm{adU}) \mathrm{K} _ {1} & \equiv \mathrm{Ue} ^ {\mathrm{U} + \mathrm{V}} e ^ {- \mathrm{U}} - e ^ {\mathrm{U} + \mathrm{V}} e ^ {- \mathrm{U}} \mathrm{U} \equiv (\mathrm{Ue} ^ {\mathrm{U} + \mathrm{V}} - e ^ {\mathrm{U} + \mathrm{V}} \mathrm{U}) e ^ {- \mathrm{U}} \\ & \equiv (e ^ {\mathrm{U}} \mathrm{V} - \mathrm{Ve} ^ {\mathrm{U}}) e ^ {- \mathrm{U}} & \text {by (23)} \\ & \equiv e ^ {\mathrm{U}} \mathrm{Ve} ^ {- \mathrm{U}} - \mathrm{V} \\ & \equiv \sum_ {n \geqslant 1} \frac {1}{n !} (\mathrm{adU}) ^ {n} \mathrm{V} & \text {by (22)} \\ & \equiv (\mathrm{adU}) \Bigl (\sum_ {n \geqslant 0} \frac {(\mathrm{adU}) ^ {n}}{(n + 1) !} \mathrm{V} \Bigr). \end{array}
\]

于是只需应用 § 2 第 11 号的注。

